% Generated draft: requires mathematical and source review.
\documentclass[11pt]{article}
\usepackage[margin=25mm]{geometry}
\usepackage{fontspec}
\setmainfont{Latin Modern Roman}
\usepackage{amsmath,amssymb,mathtools,unicode-math}
\setmathfont{Latin Modern Math}
\usepackage{enumitem,needspace,xurl,hyperref}
\setlength{\parindent}{0pt}
\setlength{\parskip}{5pt}
\setlength{\emergencystretch}{3em}
\newcommand{\N}{\mathbb{N}}
\newcommand{\Z}{\mathbb{Z}}
\newcommand{\Q}{\mathbb{Q}}
\newcommand{\R}{\mathbb{R}}
\newcommand{\C}{\mathbb{C}}
\newcommand{\F}{\mathbb{F}}
\newcommand{\A}{\mathbb{A}}
\newcommand{\PP}{\mathbb{P}}
\newcommand{\E}{\mathbb{E}}
\newcommand{\eps}{\varepsilon}
\newcommand{\dd}{\,\mathrm{d}}
\begin{document}
\section*{UnsolvedMath ID: 10400227}
\textbf{Source problem code:} catalogue entry
\subsection{Definitions and mathematical statement}
Retain \(C_2\) as the group named in the source; its definition is not provided in the statement. For a group \(G\), define its derived series by \(G^{(0)}=G\) and \(G^{(m+1)}=[G^{(m)},G^{(m)}]\). It is solvable if \(G^{(m)}=\{1\}\) for some finite \(m\). Write \(F_r\) for the free group of rank \(r\).

\emph{Questions (T. Stanford).}
\begin{enumerate}
\item Is \(C_2\) solvable; that is, does there exist a finite \(m\) with \(C_2^{(m)}=\{1\}\)?
\item Does \(C_2\) contain a free group? Literally this asks for an embedding \(F_r\hookrightarrow C_2\) for some allowed rank \(r\). If the intended meaning is a nonabelian free subgroup, the usual precise version is \(F_2\hookrightarrow C_2\). The source does not state a rank or the word ``nonabelian,'' so the distinction from an infinite cyclic subgroup must remain explicit.
\end{enumerate}
The supplied research summary identifies \(C_2\) as the second lower-central-series term of a pure braid group. Under the convention \(\gamma_1(G)=G\), \(\gamma_{j+1}(G)=[G,\gamma_j(G)]\), this would mean \(C_2=\gamma_2(P_n)=[P_n,P_n]\). However, neither this definition nor the braid index \(n\), a stabilization or direct-limit convention, or any quotient appears in the quoted statement. The editorial identification is therefore recorded conditionally and does not fix the intended group. These are two questions about that group, not an asserted equivalence between nonsolvability and existence of a free subgroup.
\subsection{Short English statement}
Stanford asks whether the group \(C_2\) is solvable and whether it contains a free subgroup. Its definition and the intended free-group rank are missing; the summary's pure-braid interpretation does not specify a braid index or other necessary conventions.
\subsection{Sources}
\begin{itemize}
\item Local UnsolvedMath registry, permanent ID 10400227, source code catalogue entry; attribution T. Stanford. The stored title is truncated after ``Problem 12.22 --- (T.''.
\item Ohtsuki, \emph{Problems on invariants of knots and 3-manifolds} (2002), Problem 12.22, PDF page 169. Supplied URL: \url{https://msp.org/gtm/2002/04/gtm-2002-04-024s.pdf}; recorded access 2026-07-29; extraction pdf-font-structure. The collection distinguishes 231 numbered bold problem, conjecture, and question headings.
\item The supplied summary interprets \(C_2\) as the second lower-central-series term of a pure braid group, but supplies no supporting definition or reference beyond the source item. Status is open and OPEN-TRIAGE. Current primary-literature status and redistribution terms for the publicly accessible author/publisher copy remain NEEDS\_REVIEW.
\end{itemize}
\end{document}
